%%%%%%%%%%%%%%%%%%%%%%%%%%%%%%%%%%%%%%%%%%%%%%%%%%%%%%%%%%%%%%%%%%%%%%
% Lớp 8 | Chương X. Một số hình khối trong thực tiễn | Hình chóp tam giác đều
% Dạng 8H0?38-1: Nhận biết các yếu tố của hình chóp tam giác đều
%%%%%%%%%%%%%%%%%%%%%%%%%%%%%%%%%%%%%%%%%%%%%%%%%%%%%%%%%%%%%%%%%%%%%%


% nguồn: De_AI_Toan8_cuoi_hoc_ki_i_1 (nhập 21/11/2026)
\begin{bt}%[8H0N38-1]%[ma:ae8c015d]
	Cho hình chóp tam giác đều $S.ABC$.
a) Hãy kể tên đỉnh, mặt đáy, các mặt bên, các cạnh bên và các cạnh đáy của hình chóp.
b) Hình chóp tam giác đều có tất cả bao nhiêu đỉnh, bao nhiêu cạnh và bao nhiêu mặt?
	\loigiai{
		a) Đối với hình chóp tam giác đều $S.ABC$: \\
- Đỉnh: $S$. \\
- Mặt đáy: tam giác đều $ABC$. \\
- Các mặt bên: $SAB$, $SBC$, $SCA$ (là ba tam giác cân bằng nhau có chung đỉnh $S$). \\
- Các cạnh bên: $SA$, $SB$, $SC$. \\
- Các cạnh đáy: $AB$, $BC$, $CA$. \\
b) Hình chóp tam giác đều có $4$ đỉnh, $6$ cạnh và $4$ mặt (gồm $1$ mặt đáy và $3$ mặt bên).
	}
\end{bt}

% nguồn: De_AI_Toan8_cuoi_hoc_ki_i_1 (nhập 21/11/2026)
\begin{bt}%[8H0N38-1]%[ma:f18ea344]
	Cho hình chóp tứ giác đều $S.ABCD$.
a) Hãy kể tên mặt đáy, các mặt bên, các cạnh bên và các cạnh đáy của hình chóp.
b) Hình chóp tứ giác đều có tất cả bao nhiêu đỉnh, bao nhiêu cạnh và bao nhiêu mặt?
	\loigiai{
		a) Đối với hình chóp tứ giác đều $S.ABCD$: \\
- Mặt đáy: hình vuông $ABCD$. \\
- Các mặt bên: $SAB$, $SBC$, $SCD$, $SDA$ (là bốn tam giác cân bằng nhau có chung đỉnh $S$). \\
- Các cạnh bên: $SA$, $SB$, $SC$, $SD$. \\
- Các cạnh đáy: $AB$, $BC$, $CD$, $DA$. \\
b) Hình chóp tứ giác đều có $5$ đỉnh, $8$ cạnh và $5$ mặt (gồm $1$ mặt đáy và $4$ mặt bên).
	}
\end{bt}

% nguồn: De_AI_Toan8_cuoi_hoc_ki_i_1 (nhập 21/11/2026)
\begin{bt}%[8H0N38-1]%[ma:26d51cb8]
	Cho hình chóp tam giác đều $S.ABC$ có cạnh đáy $AB = 6\text{ cm}$ và cạnh bên $SA = 10\text{ cm}$.
a) Nêu độ dài các cạnh đáy và các cạnh bên còn lại của hình chóp.
b) Tính chu vi mặt đáy và tổng độ dài tất cả các cạnh của hình chóp.
	\loigiai{
		a) Vì $S.ABC$ là hình chóp tam giác đều nên đáy $ABC$ là tam giác đều và các cạnh bên bằng nhau: \\
$BC = CA = AB = 6\text{ cm}$ và $SB = SC = SA = 10\text{ cm}$. \\
b) Chu vi mặt đáy $ABC$ là $3 \cdot 6 = 18\text{ cm}$. \\
Tổng độ dài ba cạnh bên là $3 \cdot 10 = 30\text{ cm}$. \\
Tổng độ dài tất cả các cạnh của hình chóp là $18 + 30 = 48\text{ cm}$.
	}
\end{bt}

% nguồn: De_AI_Toan8_cuoi_hoc_ki_i_1 (nhập 21/11/2026)
\begin{bt}%[8H0N38-1]%[ma:fbecc37c]
	Cho hình chóp tứ giác đều $S.ABCD$ có cạnh đáy $AB = 8\text{ cm}$ và cạnh bên $SA = 12\text{ cm}$.
a) Nêu độ dài các cạnh đáy và các cạnh bên còn lại của hình chóp.
b) Tính chu vi mặt đáy và tổng độ dài tất cả các cạnh của hình chóp.
	\loigiai{
		a) Vì $S.ABCD$ là hình chóp tứ giác đều nên mặt đáy $ABCD$ là hình vuông và bốn cạnh bên bằng nhau: \\
$BC = CD = DA = AB = 8\text{ cm}$ và $SB = SC = SD = SA = 12\text{ cm}$. \\
b) Chu vi mặt đáy $ABCD$ là $4 \cdot 8 = 32\text{ cm}$. \\
Tổng độ dài bốn cạnh bên là $4 \cdot 12 = 48\text{ cm}$. \\
Tổng độ dài tất cả các cạnh của hình chóp là $32 + 48 = 80\text{ cm}$.
	}
\end{bt}

% nguồn: De_AI_Toan8_cuoi_hoc_ki_i_1 (nhập 21/11/2026)
\begin{bt}%[8H0N38-1]%[ma:d4f31c1b]
	So sánh các đặc điểm sau giữa hình chóp tam giác đều và hình chóp tứ giác đều:
a) Hình dạng mặt đáy và hình dạng các mặt bên;
b) Vị trí chân đường cao hạ từ đỉnh chóp xuống mặt đáy.
	\loigiai{
		a) Về hình dạng mặt đáy và mặt bên: \\
- Hình chóp tam giác đều có mặt đáy là tam giác đều và $3$ mặt bên là các tam giác cân bằng nhau. \\
- Hình chóp tứ giác đều có mặt đáy là hình vuông và $4$ mặt bên là các tam giác cân bằng nhau. \\
b) Về chân đường cao hạ từ đỉnh chóp xuống mặt đáy: \\
- Ở hình chóp tam giác đều, chân đường cao là trọng tâm (đồng thời là tâm đường tròn ngoại tiếp) của tam giác đều đáy. \\
- Ở hình chóp tứ giác đều, chân đường cao là giao điểm hai đường chéo (tâm) của hình vuông đáy.
	}
\end{bt}

% nguồn: De_AI_Toan8_cuoi_hoc_ki_i_1 (nhập 21/11/2026)
\begin{bt}%[8H0N38-1]%[ma:acc35115]
	Cho hình chóp tam giác đều $S.MNP$. Gọi $O$ là trọng tâm của tam giác $MNP$ và $H$ là trung điểm của cạnh $NP$.
a) Nêu tên đường cao và một trung đoạn của hình chóp $S.MNP$.
b) So sánh độ dài các đoạn thẳng $SM$, $SN$, $SP$ và cho biết tam giác $SNP$ là tam giác gì.
	\loigiai{
		a) Vì $O$ là trọng tâm của tam giác đều $MNP$ nên $SO$ là đường cao của hình chóp $S.MNP$. \\
Vì $H$ là trung điểm của cạnh đáy $NP$ nên $SH$ là một trung đoạn của hình chóp $S.MNP$ (và $SH \perp NP$). \\
b) Các cạnh bên của hình chóp tam giác đều bằng nhau nên $SM = SN = SP$. \\
Do $SN = SP$ nên tam giác $SNP$ là tam giác cân tại $S$.
	}
\end{bt}

% nguồn: De_AI_Toan8_cuoi_hoc_ki_i_1 (nhập 21/11/2026)
\begin{bt}%[8H0N38-1]%[ma:eff758d0]
	Cho hình chóp tứ giác đều $S.ABCD$. Gọi $O$ là giao điểm hai đường chéo $AC$ và $BD$, gọi $K$ là trung điểm của cạnh $CD$.
a) Nêu tên đường cao và một trung đoạn của hình chóp $S.ABCD$.
b) Chứng tỏ $OK \perp CD$ và tính $OK$ khi biết cạnh đáy $AB = 10\text{ cm}$.
	\loigiai{
		a) Vì $O$ là giao điểm hai đường chéo của hình vuông đáy $ABCD$ nên $SO$ là đường cao của hình chóp $S.ABCD$. \\
Vì $K$ là trung điểm của cạnh đáy $CD$ nên $SK$ là một trung đoạn của hình chóp $S.ABCD$. \\
b) Trong tam giác $ACD$, $O$ là trung điểm $AC$ và $K$ là trung điểm $CD$, nên $OK$ là đường trung bình của tam giác $ACD$. \\
Suy ra $OK \parallel AD$ và $OK = \dfrac{1}{2}AD$. Vì $AD \perp CD$ nên $OK \perp CD$. \\
Khi $AB = 10\text{ cm}$ thì $AD = 10\text{ cm}$, do đó $OK = \dfrac{10}{2} = 5\text{ cm}$.
	}
\end{bt}

% nguồn: De_AI_Toan8_cuoi_hoc_ki_i_1 (nhập 21/11/2026)
\begin{bt}%[8H0N38-1]%[ma:92c5c788]
	Một hình chóp tam giác đều có tất cả các cạnh đều bằng nhau và bằng $5\text{ cm}$.
a) Các mặt bên và mặt đáy của hình chóp là hình gì?
b) Tính tổng độ dài tất cả các cạnh của hình chóp đó.
	\loigiai{
		a) Vì tất cả các cạnh đáy và cạnh bên đều bằng $5\text{ cm}$ nên mặt đáy và cả $3$ mặt bên đều là các tam giác đều có cạnh bằng $5\text{ cm}$. \\
b) Hình chóp tam giác đều có $6$ cạnh bằng nhau nên tổng độ dài tất cả các cạnh là $6 \cdot 5 = 30\text{ cm}$.
	}
\end{bt}
